\documentclass[12pt]{article}
% Préambule commun à tous les documents extraits de « La Revue CAPES »
\usepackage[T1]{fontenc}
\usepackage[utf8]{inputenc}
\usepackage{lmodern}
\usepackage{amsmath,amssymb,amsthm,amsfonts,mathrsfs}
\usepackage{setspace}
\usepackage{listings}
\usepackage{pgfplots}
\pgfplotsset{compat=1.18}
\usepackage{longtable}
\usepackage{adjustbox}
\usepackage{lettrine}
\usepackage{tkz-tab}
\usepackage{changepage}
\usepackage{pdflscape}
\usepackage{graphicx}
\usepackage{tikz}
\usepackage{xcolor}
\usepackage[a4paper,top=2cm,bottom=2.5cm,left=2.5cm,right=2cm]{geometry}
\usepackage{fancyhdr}
\usepackage{draftwatermark}
\SetWatermarkText{La RevueCapes}
\SetWatermarkLightness{0.9}
\SetWatermarkScale{2}
\usepackage{hyperref}
\hypersetup{colorlinks=true,linkcolor=blue!50!black,urlcolor=blue!50!black}

\definecolor{revue}{RGB}{20,60,120}

\pagestyle{fancy}
\fancyhf{}
\renewcommand{\headrulewidth}{0.4pt}
\setlength{\headheight}{15pt}

% \entete{Matière}{Titre}{Type de document}
\newcommand{\entete}[3]{%
  \fancyhead[L]{\small\textsc{La Revue CAPES} -- #1}%
  \fancyhead[R]{\small #2}%
  \fancyfoot[C]{\thepage}%
  \thispagestyle{plain}%
  \begin{center}
    {\color{revue}\rule{\textwidth}{1.2pt}}\\[4pt]
    {\small\textsc{Concours directs CAP-CEG / CAPES -- Burkina Faso}}\\[6pt]
    {\LARGE\bfseries\color{revue} #2}\\[6pt]
    {\large\itshape #1 -- #3}\\[2pt]
    {\color{revue}\rule{\textwidth}{1.2pt}}
  \end{center}
  \vspace{0.5cm}%
}

% Encadré pour les remarques ajoutées dans les corrigés
\newcommand{\remarque}[1]{\par\medskip\noindent\fcolorbox{revue}{revue!6}{\parbox{\dimexpr\linewidth-2\fboxsep-2\fboxrule}{\textbf{\color{revue}Remarque.} #1}}\par\medskip}


\begin{document}
\entete{Mathématiques}{Sujet type n°14}{Énoncé}
\begin{spacing}{1.5}

\textbf{Exercice 1}

Le plan est rapporté à un repère orthonormal. 

1. On considère une courbe $C$ du plan définie comme l'ensemble des points $M$ de coordonnées $x(t)$ et $y(t)$, dépendant d'un paramètre réel $t$. 

On donne, pour $t\in \mathbb{R} $: 

$x(t)=t+\frac{t^2}{2}$ et $y(t)=t-\frac{t^2}{2}$

a. Étudier, dans le même tableau de variations, les variations de $x(t)$ et $y(t)$ en fonction de $t$. 

b. Étudier les points de la courbe C en lesquels la tangente est parallèle à l'un des axes du repère. 

c. Préciser les points d'intersection de $C$ avec les axes $Ox$ et $Oy$. Donner un vecteur 
directeur des tangentes en ces points. 

d. Donner graphiquement l'allure de la courbe $C$, en positionnant les points remarquables 
déterminés. 

e. Donner l'équation cartésienne $v(x, y) = 0$ de la courbe $C$. 

2. On se propose de déterminer précisément ce qu'est la courbe $C$. 

a. Le plan étant assimilé au plan complexe, on considère l'application h qui, à tout point $M$ 
d'affixe $z$ du plan, associe un point $M'$ d'affixe $Z$, définie par : 

$Z = \frac{1}{\sqrt{2}}(1+i)z$

 Donner de façon précise la nature de l'application h. 

b. Soit $M$  le point d'affixe $z(t) = x(t) + iy(t)$. 

Déterminer les coordonnées $X(t)$ et $Y(t)$ de $M'$, image de $M$ par $h$. 

Donner une équation cartésienne $\varphi (X,Y)=0$ de la courbe $C'$ des points $M'$. 

Comment appelle-t-on la courbe $C'$ ? 

En déduire le nom de la courbe $C$. 


\medskip\textbf{Exercice 2}

 Donner la probabilité, l'espérance mathématique et la variance de la variables aléatoires $X$ lorsque $X$ suit chacune des lois discrètes usuelles.
 
\medskip\textbf{Exercice 3}

M.Barro a été sélectionné pour participé à un jeu télévisé. Chaque mois on lui verse, pendant une durée maximale de 12 mois,une somme d'argent dont le montant initial est 600 000 FCFA le premier mois.

 Le versement augmente, sur le modèle des intérêts composés avec un taux d'intérêt $i=0.03$ chaque mois. Ainsi, le $2^{e}$ mois M.Barro touchera 618 000 fcfa, etc. 
 
 On désigne par $U_n$ le montant, exprimé en fcfa, versé à M.Barro le $n^{ieme}$ mois. Ainsi, $U_1$=600 000 fcfa.
 
 On arrondira éventuellement les résultats à l'unité.

1. Calculer la somme exacte que M.Barro touchera le $3^{e}$ mois.

2. Quelle est la nature de la suite $(U_n)$? 

Exprimer $U_n$ en fonction de $ n$.

3. Calculer la somme que M.Barro touchera le dernier mois.

4. M.Barro veut calculer la somme totale qu'il a gagnée durant ces 12 mois. On pose $S=U_1+U_2+...+
 U_{12}$.

Déterminer la somme $S$.


\quad
\quad

\end{spacing}
\end{document}
